\documentclass[11pt,a4paper]{article}
\usepackage[utf8]{inputenc}
\usepackage[T1]{fontenc}
\usepackage{lmodern}
\usepackage[french]{babel}
\usepackage{amsmath,amssymb,mathtools}
\usepackage{icomma}
\usepackage{xcolor}
\usepackage{tikz}
\usepackage{pgfplots}
\usepackage[most]{tcolorbox}
\usepackage{enumitem}
\usepackage{fancyhdr}
\usepackage[hidelinks]{hyperref}
\usepackage{titlesec}
\usepackage[a4paper,margin=2.2cm,headheight=26pt,headsep=10pt]{geometry}
\usetikzlibrary{arrows.meta,positioning,calc,intersections,angles,quotes}
\usepgfplotslibrary{fillbetween}
\pgfplotsset{compat=1.18}
\definecolor{primary}{RGB}{30,75,145}
\definecolor{primarydark}{RGB}{20,50,100}
\definecolor{defcol}{RGB}{0,114,178}
\definecolor{thmcol}{RGB}{0,135,110}
\definecolor{excol}{RGB}{0,130,85}
\definecolor{remarkcol}{RGB}{120,70,180}
\definecolor{attncol}{RGB}{200,40,55}
\definecolor{methcol}{RGB}{85,85,100}
\definecolor{areacol}{RGB}{120,160,230}
\definecolor{areacol2}{RGB}{230,150,80}
\newcommand{\dd}{\mathrm{d}}
\newcommand{\R}{\mathbb{R}}
\newcommand{\ua}{\,\text{u.a.}}
\newcommand{\tg}{\tan}\newcommand{\cotg}{\cot}
\tcbset{boxbase/.style={enhanced,breakable,boxrule=0.9pt,arc=2.5pt,left=3mm,right=3mm,top=2.3mm,bottom=2.3mm,fonttitle=\bfseries,coltitle=white,toptitle=0.6mm,bottomtitle=0.6mm}}
\newtcolorbox{methode}[1][]{boxbase,colback=methcol!7,colframe=methcol,sharp corners,title={\textsc{Comment faire}\ifx\relax#1\relax\relax\else\textmd{ : \itshape #1}\fi}}
\newtcolorbox{exemple}[1][]{boxbase,colback=excol!5,colframe=excol,title={\textsc{Exemple}\ifx\relax#1\relax\relax\else\textmd{ : \itshape #1}\fi}}
\newtcolorbox{idee}[1][]{boxbase,colback=defcol!5,colframe=defcol,title={\textsc{Idée clé}\ifx\relax#1\relax\relax\else\textmd{ : \itshape #1}\fi}}
\newtcolorbox{remarque}[1][]{boxbase,colback=remarkcol!6,colframe=remarkcol,title={\textsc{Remarque}\ifx\relax#1\relax\relax\else\textmd{ : \itshape #1}\fi}}
\newtcolorbox{attention}[1][]{boxbase,colback=attncol!6,colframe=attncol,title={\textsc{Attention}\ifx\relax#1\relax\relax\else\textmd{ : \itshape #1}\fi}}
\newtcolorbox{exo}[1][]{boxbase,colback={primary!4},colframe=primary,title={\textsc{Exercice #1}}}
\newtcolorbox{corr}[1][]{boxbase,colback={thmcol!5},colframe=thmcol,title={\textsc{Corrigé #1}}}
\tikzset{axe/.style={->,>=Stealth,black!65,line width=0.7pt},courbe/.style={primary,line width=1.3pt},courbe2/.style={areacol2!90!black,line width=1.3pt},dvert/.style={gray,dashed,line width=0.7pt}}
\pagestyle{fancy}\fancyhf{}
\fancyhead[L]{\small\textcolor{primarydark}{\textbf{Fiche 7} — Intégrale d'une fonction continue}}
\fancyhead[R]{\small\textcolor{primarydark}{\textsc{Thème 4 — Intégrale définie}}}
\fancyfoot[C]{\small\thepage}
\renewcommand{\headrulewidth}{0.8pt}
\titleformat{\section}{\large\bfseries\color{primarydark}}{\thesection}{0.6em}{}

\begin{document}

\section*{Niveau Moyen — Corrigé détaillé}

\begin{corr}[1 — composée trigonométrique]
On fait apparaître le facteur $2$ demandé par le $\dd U$ ($U=2x$, $\dd U=2\,\dd x$) :
\[
A=\int_0^{\pi/6}\sin(2x)\,\dd x=\frac12\int_0^{\pi/6}\sin(2x)\cdot 2\,\dd x
=\frac12\Bigl[-\cos(2x)\Bigr]_0^{\pi/6}.
\]
Or $\cos\bigl(2\cdot\tfrac{\pi}{6}\bigr)=\cos\tfrac{\pi}{3}=\tfrac12$ et $\cos 0=1$, donc
\[
A=\frac12\Bigl(-\cos\tfrac{\pi}{3}+\cos 0\Bigr)=\frac12\Bigl(-\frac12+1\Bigr)=\frac12\cdot\frac12=\boxed{\,\frac14\,}\ua.
\]
\end{corr}

\begin{corr}[2 — exponentielle composée]
Une primitive de $e^{2x}$ est $\tfrac12 e^{2x}$ (on compense le facteur $2$). Donc
\[
B=\left[\frac12 e^{2x}\right]_0^{1/2}=\frac12 e^{2\cdot 1/2}-\frac12 e^{0}=\frac12 e^{1}-\frac12
=\frac{e-1}{2}\approx\frac{1{,}718}{2}\approx\boxed{\,0{,}86\,}\ua.
\]
\textbf{Vérification} : $\bigl(\tfrac12 e^{2x}\bigr)'=\tfrac12\cdot 2\,e^{2x}=e^{2x}$. \checkmark
\end{corr}

\begin{corr}[3 — le $\dd U$ presque présent]
Posons $U=x^2+1$, donc $\dd U=2x\,\dd x$, soit $x\,\dd x=\tfrac12\,\dd U$. On adapte les bornes :
$x=0\Rightarrow U=1$ et $x=1\Rightarrow U=2$. Ainsi
\[
C=\int_0^1 x\,e^{x^2+1}\,\dd x=\frac12\int_1^2 e^{U}\,\dd U=\frac12\bigl[e^{U}\bigr]_1^2
=\frac12\bigl(e^2-e\bigr)=\boxed{\,\frac{e^2-e}{2}\,}\approx 2{,}34\ua.
\]
\textbf{Astuce} : sans changer les bornes, on écrit directement
$\tfrac12\bigl[e^{x^2+1}\bigr]_0^1=\tfrac12(e^2-e)$ — même résultat.
\end{corr}

\begin{corr}[4 — relation de Chasles et inversion]
\textbf{a)} La relation de Chasles donne $\displaystyle\int_1^2 f+\int_2^4 f=\int_1^4 f$, d'où
\[
\int_2^4 f(x)\,\dd x=\int_1^4 f-\int_1^2 f=7-2=\boxed{\,5\,}.
\]
\textbf{b)} Inverser les bornes change le signe :
\[
\int_4^2 f(x)\,\dd x=-\int_2^4 f(x)\,\dd x=-5=\boxed{\,-5\,}.
\]
\end{corr}

\begin{corr}[5 — linéarité pour simplifier]
Par linéarité, $D=\displaystyle\int_0^{\pi/2}\cos x\,\dd x+2\int_0^{\pi/2}\sin x\,\dd x$. Avec les
primitives $\sin x$ et $-\cos x$ :
\[
D=\bigl[\sin x\bigr]_0^{\pi/2}+2\bigl[-\cos x\bigr]_0^{\pi/2}
=\bigl(1-0\bigr)+2\bigl(-\cos\tfrac{\pi}{2}+\cos 0\bigr)=1+2\bigl(0+1\bigr)=\boxed{\,3\,}\ua.
\]
On peut aussi regrouper d'emblée : $D=\bigl[\sin x-2\cos x\bigr]_0^{\pi/2}=(1-0)-(0-2)=3$.
\end{corr}

\end{document}
